\exercisesection{2}

\begin{enumerate}
\def\labelenumi{\arabic{enumi}.}
\item
  设 \(\Omega\) 是域 K 的一个扩域，E、F 是包含于 \(\Omega\) 中的 K 的两个扩域，且在 K 上线性不交，L 是包含于 E 中的 K 的扩域。令 \(f\) 为 E 到 L 的位，它在 K 上化为恒等自同构。证明存在唯一的位，定义在 K(E \(\cup\) F) 上，取值于 K(L \(\cup\) F)，延拓 \(f\)，并在 F 上化为恒等自同构。（若 A 是 \(f\) 的环，注意 \(f\) 可唯一延拓为同态 \(g\)，从 F{[}A{]} 到 K(L \(\cup\) F)，在 F 上化为恒等，并且若 m 是 \(f\) 的理想，则每个满足 \(z \in\) F{[}A{]} 且 \(g(z) = 0\) 的元素 z 都可写成 \(xu\)，其中 \(x \in m\)，\(u \in\) F{[}A{]} 且 \(g(u) \neq 0\)。最后注意 K(E \(\cup\) F) 是 F{[}A{]} 的分式域。）
\item
  设 K、K' 是域 k 的两个扩域，f 和 \(f'\) 分别是 K 和 \(K'\) 到同一代数闭域 L 的位。假设 f 与 \(f'\) 在 k 上的限制相同。证明存在 K 与 \(K'\) 的复合扩域 (F, i, i')（《代数》第八章 § 8）以及 F 到 L 的位 g，使得 \(f(x) = g(i(x))\) 对 \(x \in \mathbf{K}\) 成立，且 \(f'(x) = g(i'(x'))\) 对 \(x' \in \mathbf{K}'\) 成立。（令 V、\(V'\) 分别为 f、\(f'\) 的环，A 为它们与 k 的公共交，且 h: \(V \otimes_{A} V' \to L\) 为满足 \(h(a \otimes a') = f(a)f'(a')\)（其中 \(a \in V\)、\(a' \in V'\)）的同态。利用 V 和 \(V'\) 是平坦 A-模这一事实（§ 3 第 5 节引理 1），证明若 \(a \neq 0\)、\(a' \neq 0\)，则 \(a \otimes a' = (a \otimes 1)(1 \otimes a')\) 不是零因子；若 S 是 B = V \(\otimes_{A} V'\) 中由这些元素组成的乘法子集，
\end{enumerate}

\(S^{-1}B\) 含有子域 \(K_{1}, K_{1}'\)，它们分别同构于 K 和 \(K'\)，并且它是由 \(K_{1} \cup K_{1}'\) 生成的环。证明若 q 是 h 的核，则 B 中存在素理想 p，使 \(p \subset q\) 且 \(p \cap S = \varnothing\)，并证明 F 可取为 \(S^{-1}B/S^{-1}p\) 的分式域。

\begin{enumerate}
\def\labelenumi{\arabic{enumi}.}
\setcounter{enumi}{2}
\tightlist
\item
  \begin{enumerate}
  \def\labelenumii{(\alph{enumii})}
  \tightlist
  \item
    设 \(\mathbf{K}\) 为域，\(f\) 是 \(\mathbf{K}\) 上取值于可序域 \(\mathbf{L}\) 的位（《代数》第六章 § 2 习题 8）；证明 \(\mathbf{K}\) 可赋序。（注意，对于元素有限族 \((a_i)_{1\leq i\leq n}\) 属于 \(\mathbf{K}\)，总存在指标 \(j\)，使 \(f(a_i / a_j)\neq \infty\) 对所有 \(i\) 成立。）
  \end{enumerate}
\end{enumerate}

若 \((x_{i})_{1\leqslant i\leqslant n}\) 是 K 中元素的序列，且 \(f\left(\sum_{i = 1}^{n}x_{i}^{2}\right)\neq \infty\)，证明 \(f(x_{i})\neq \infty\) 对所有 \(i\) 成立。

\begin{enumerate}
\def\labelenumi{(\alph{enumi})}
\setcounter{enumi}{1}
\item
  设 K 为有序域，G 为 K 的子群；证明环 \(F(G)\)（由 K 中相对于 G 不是无限大的元素组成；同上，习题 11）是 K 的赋值环，而理想 \(I(G)\)（由 \(F(G)\) 中相对于 G 无穷小的元素组成）是此环的极大理想；因而 K 上存在一个取值于 \(k(G) = F(G)/I(G)\) 的位，它在 G 上化为恒等；称此位为 K 关于 G 的典范位。
\item
  证明若 K 中不存在异于 G 且与 G 可比较的 G 的扩域，则 \(k(G)\) 在 G 上代数（注意，若 \(t \in K\) 不属于 G，则域 \(G(t)\) 含有相对于 G 无穷小的元素 \(u \neq 0\)，且 \(G(t)\) 是 \(G(u)\) 的代数扩张）。
\end{enumerate}

¶ 4. 若有序域 K 对于 K 中所有 \(x \geqslant 0\) 都存在 \(y \in \mathbf{K}\) 使 \(x = y^2\)，则称 K 为欧氏的。

\begin{enumerate}
\def\labelenumi{(\alph{enumi})}
\item
  设 K 为欧氏有序域，G 为 K 的子域，且 K 中不存在异于 G 且与 G 可比较的 G 的扩域。证明：若 f 是 K 到有序域 L 的位，L 是 G 的代数扩域，且 f 在 G 上化为恒等，则 f 等价于 K 关于 G 的典范位。（可设 f 取值于 L 的一个极大有序代数扩域 N。注意 G 是欧氏的，并且若 \(x \in G\)、\(y \in K\) 且 x \textless{} y，则 \(f(y) = \infty\) 或 \(f(y) > x\)。若 A 和 p 是 f 的环和理想，依次证明 \(p \subset I(G)\)、\(F(G) \subset A\)、\(I(G) \subset p\) 以及 \(A \subset F(G)\)，并注意 N 与 G 可比较。）
\item
  设 K 为欧氏有序域，f 为一个取值于极大有序域 L 的 K 上的位，A 和 p 为 f 的环和理想。~令 G 为满足 \(E \cap p = 0\) 的 K 的子域 E 中的极大子域。证明 \(G \subset A\)，且 G 在 K 中代数闭，因而是欧氏的。证明 K 中不存在异于 G 且与 G 可比较的 G 的扩域。此外，L 中由 \(f(A)\) 生成的子域在 \(f(G)\) 上代数，因而 f 等价于 K 关于 G 的典范位。
\item
  ¶ 6. 设 K 为极大有序域，L 为 K 的异于 K 的极大有序子域，且 K 与 L 可比较（《代数》第六章 § 2 习题 15）。证明不存在取值于 L、在 L 上化为恒等自同构的 K 的位 f。
\end{enumerate}

¶ 5. 设 K 为域，f 为一个取值于极大有序域 L 的 K 上的位。

\begin{enumerate}
\def\labelenumi{(\alph{enumi})}
\item
  对所有 \(\alpha \in \mathbf{K}\)，证明 \(f\) 存在一个延拓，要么延拓到 \(\mathbf{K}(\sqrt{\alpha})\)，要么延拓到 \(\mathbf{K}(\sqrt{-\alpha})\)，并且该延拓是取值于 \(\mathbf{L}\) 的位。（根据以下两种情形之一进行讨论：存在 \(a \in \mathbf{K}\) 使 \(f(a^2\alpha)\) 既非 0 也非 \(\infty\)，或者 \(f(x^2\alpha)\) 等于 0 或 \(\infty\) 对所有 \(x \in \mathbf{K}\) 成立。）
\item
由 (a) 推出，存在 K 的一个扩域 E，它是极大有序域，并且 f 到 E 的延拓是取值于 L 的 E 上的位。（将其归结为 K 为欧氏的情形，并使用习题 4 (a)。）
\end{enumerate}

\begin{enumerate}
\def\labelenumi{\arabic{enumi}.}
\setcounter{enumi}{5}
\tightlist
\item
  设 A 为整闭整环，K 为其分式域，p 为 A 的素理想，k 为 A/p 的分式域。~证明若 \(A_{p}\) 不是赋值环，则 K 中存在支配 \(A_{p}\) 的赋值环 V，且其剩余域是 k 的超越扩张（使用《代数》第五章 § 2 习题 14(b)）。
\end{enumerate}

